\section{Introduction}
\label{sec:intro}

The public ALMA archive holds thousands of hours of millimetre and
submillimetre observations taken for other reasons, mostly studies of dust
discs and planet formation. Their visibilities allow a spectrum at the
position of the star at the centre of each field, at a resolution and depth
set by the original observer's correlator setup. This paper searches those
spectra for a narrowband transmitter at the star's own position. Nothing of
the kind survived.

Searches for \emph{technosignatures} --- detectable imprints of technology on
the electromagnetic spectrum, intended or not
\citep{Sheikh2020,SETIreview2025} --- have examined many thousands of stars
over six decades, yet sample only \HaystackFrac{} of the space of transmitter
frequency, bandwidth, power and duty cycle \citep{Wright2018}, and
essentially all of that effort lies below ${\sim}30$\,GHz. No published radio
technosignature search above ${\sim}30$\,GHz has reached a star within
40\,pc: the one published search in that range worked in ALMA Band~3 toward
targets at \MasonNearKpc\,kpc and beyond \citep{Mason2025}, and few searches
at any frequency reprocess existing interferometric archives
\citep{White2026}. The band is attractive because a carrier arrives less
distorted and a given aperture delivers more gain per watt supplied, and hard
because archived correlator channels are thousands of times coarser than a
dedicated narrowband backend and are crowded with molecular lines, a dense
background of real astrophysical signals that a threshold search must
separate from an artificial one
(\S\ref{sec:background},~\S\ref{sec:mmdifferent}).

The signal sought is an \emph{unresolved spectral carrier}, emission confined
to a few channels, unresolved at the stellar position, and persistent enough
to recur between epochs. Each qualifier is a restriction. A carrier broader
than the channel width would not be recovered, nor one that fills the
synthesised beam, nor one that appears once and never again --- and it is
recurrence rather than brightness that separates an interesting event from a
detection \citep{Sheikh2021}. The frequency is allowed to drift, as it must
for a transmitter carried on a rotating, orbiting body, and the drift
searched is bounded by a line-of-sight acceleration rather than by a
frequency (\S\ref{sec:statistic}). The channels used are
\ChanLoA--\ChanHiA\,kHz wide, fine only relative to this archival sample, so
``narrowband'' here describes the assumed transmitter and never the resolving
power of the data. The targets were chosen by archival disc and
planet-formation programmes and not for this search, so every statistical
statement that follows is about an archive-selected sample and not about the
solar neighbourhood (\S\ref{sec:excludes}).
